\section{Methods}

\subsection{Data acquisition}

All data were retrieved programmatically from public NASA APIs; no value in this
study was simulated, imputed or hand-entered. Plant and crew data came from the
OSDR file API at
\path{https://visualization.osdr.nasa.gov/biodata/api/v2/dataset/{OSD-ID}/files/},
preferring OSDR's curated \texttt{*\_TRANSFORMED.csv} products, whose column
names carry the unit alongside the analyte. Each downloaded file is cached and a
SHA-256 digest recorded in \path{fair/checksums.sha256}, so a re-run can be
verified byte-for-byte against ours.

The LSDA catalog was harvested from the JSON query endpoint at
\path{https://nlsp.nasa.gov/explore/api/v1/query}. Two behaviours of that
endpoint, established by direct probing rather than documentation, determined
the approach. Its \texttt{query} parameter does not filter server-side --- every
search term returns the full corpus of \LSDATotal{} records and the term affects
only relevance ordering --- so the entire catalog was harvested and filtered
locally. It also returns HTTP 502/503 under rapid request rates, so retrieval
was throttled to one request per two seconds with exponential-backoff retry.
Pagination uses a \texttt{page} parameter at 20 records per page; the
\texttt{from}, \texttt{offset}, \texttt{skip} and \texttt{start} parameters are
accepted but silently ignored, each returning page one.

\subsection{Mineral and nutrition annotation of the LSDA catalog}

Each experiment record was matched against a curated ontology of twelve minerals
and their clinical surrogates (for example, ferritin and transferrin for iron;
parathyroid hormone and bone resorption markers for calcium) and a set of
nutritional terms, searched across the measured-parameter list, research areas,
keywords, title and description. All terms are word-bounded: an unbounded
\texttt{iron} matches \emph{env\underline{iron}ment} and alone produced 1{,}266
spurious matches during development. Experiments were separately flagged for
whether a mineral appeared in the \emph{measured parameter} list, which states
what was actually assayed, as opposed to anywhere in the record. LSDA's
free-text \texttt{archive\_status} was mapped to four actionable categories:
openly available, request required, held elsewhere, and never submitted.

\subsection{Unit harmonisation}
\label{sec:units}

All elemental concentrations were converted to $\mu$g/g dry weight. OSD-745
reports mg/kg, numerically identical to $\mu$g/g, requiring no rescaling.
OSD-655 names five columns --- calcium, potassium, magnesium, phosphorus and
sulfur --- \texttt{*\_milligram\_per\_gram\_of\_dry\_weight}, and
\path{OSD-655_dictionary_DataDictionary_OSD-655.csv} repeats
\texttt{unit: milligram\_per\_gram} for each.

\begin{smacaution}[title=The values are not mg/g]
Read as stated, mean concentrations would be $2943\%$ (K), $805\%$ (Mg),
$651\%$ (Ca), $551\%$ (S) and $367\%$ (P) of dry weight. Read as $\mu$g/g they
fall in the normal range for leafy \emph{Brassica} tissue and agree with the
independently-scaled OSD-745 lettuce measurements. The six elements OSD-655
labels in $\mu$g/g (Al, B, Cu, Fe, Mn, Zn) need no correction and already agree
across studies, which localises the defect to the five mg/g labels. Because the
archived data dictionary carries the same error, the defect is in the curation
metadata and not only in the file headers.
\end{smacaution}

The correction is applied through an explicit override table in which each entry
records the stated unit, the applied unit and the evidence for the change;
\NUnitOverrides{} overrides fired. Corrections are never silent: the
harmonisation report in \path{fair/harmonisation_report.txt} lists every one.

\subsection{Quality control}

Four gates run between curation and analysis, independent of the harmonisation
code so that the two must agree. A \emph{physical plausibility} gate fails the
build if any element exceeds $10\%$ of dry weight --- the check that catches a
1000-fold unit error. A \emph{cross-study agreement} gate fails if the two
crops' ground-control means differ by more than tenfold; the threshold sits two
orders of magnitude below the error it targets, and is deliberately loose
because genuine between-species differences here are large (ground-control
sodium differs 5.4-fold between the crops). Further gates check that each crop
has both flight and ground arms, and that missingness is structural rather than
sporadic. All gates passed (\QCAllPassed).

\subsection{Compositional analysis}
\label{sec:clr}

ICP-OES reports concentrations --- element mass per unit dry tissue mass ---
which are compositional and share a fixed denominator \citep{aitchison1982}. If
flight and ground tissue differ in dry-matter composition, every element's
concentration shifts together even with no change in uptake. This is not
hypothetical here: in the raw OSD-655 data eight elements rise in concert, the
signature of a denominator effect rather than eight independent responses.

Each sample was therefore also expressed as a centred log-ratio,
$\mathrm{clr}(x_i) = \log x_i - \frac{1}{D}\sum_{j=1}^{D} \log x_j$, which
divides by the sample's geometric mean and so removes any per-sample
multiplicative factor exactly. A CLR contrast asks whether an element changed
\emph{relative to the rest of the ionome}, and is invariant to dilution.

The composition must use the same parts in every sample. Missingness here is
structural --- VEG-03A did not assay boron, so all twelve of its samples lack
that element while being otherwise complete --- so incomplete \emph{columns}
were dropped rather than incomplete rows, which would have discarded an entire
mission. Aluminium is excluded from the composition and analysed separately: it
is a non-nutrient whose passive uptake tracks substrate carryover rather than
plant physiology.

\subsection{Statistics}

Flight-versus-ground contrasts used two-sided Mann--Whitney $U$ tests, chosen
over $t$-tests because group sizes are 11--18 and normality is not established.
Effect sizes are reported as rank-biserial correlation alongside the descriptive
fold change. Multiplicity was controlled within each crop across elements by the
Benjamini--Hochberg procedure \citep{benjamini1995}; $q<0.05$ is reported as
significant.

\subsection{Classification}

Flight versus ground was classified from the ionome using $L_2$-penalised
logistic regression and a random forest, each in a pipeline whose scaler is
fitted inside the training fold. Only the \MLNFeatures{} elements complete
across all \NSamples{} samples were used (\MLFeatures); imputation was avoided
because the missingness is structural and an imputed value would be an
invention.

Two grouped cross-validation schemes were used, since samples from one mission
share a pillow batch, a fertiliser charge and a harvest date, and a random split
would place siblings across the fold boundary. \emph{Leave-one-mission-out}
tests generalisation to an unseen growout; \emph{leave-one-crop-out} trains on
one species and tests on the other. Hyperparameters were selected by an inner
grouped cross-validation on training folds only, so no reported score is one
that tuning maximised. Every model is reported against stratified and
majority-class dummy classifiers evaluated under the identical scheme: at this
sample size an accuracy figure without its baseline carries no information.
Confidence intervals are 2000-sample percentile bootstraps. There are no
fallback metrics anywhere in the pipeline; a model that fails to fit stops the
build.

\subsection{Supply and serving calculations}

Per-pillow water delivery and harvest mass were parsed from the OSD-655
supplements and joined to the elemental data on pillow identifier, which the
ICP-OES sample names carry. Two limits are propagated into every derived figure:
the supplements cover VEG-04A only, matching \SupplyMatched{} of
\SupplyTotal{} Tokyo Bekana samples; and harvest masses are \emph{fresh} while
concentrations are per \emph{dry} gram, with no dry-matter fraction archived.
Mineral content is therefore bounded across a dry-matter fraction of
$5$--$10\%$ rather than point-estimated. Daily requirements are taken verbatim
from NASA JSC-67378 Rev 1 \citep{jsc67378}, not from terrestrial dietary
reference intakes, which differ --- NASA caps iron at 8\,mg/d specifically
because of spaceflight-induced increases in iron stores.

\subsection{Reproducibility}

Python 3.11 with exactly pinned dependencies. The pipeline runs end to end via
\texttt{make all}: fetch, curate, quality control, analyse, figure generation,
manuscript build. Every numeric value in this manuscript is injected from the
analysis outputs as a \LaTeX{} macro generated by
\path{scripts_generate_numbers.py}; no statistic is typed into the prose, and a
macro referenced without a definition halts the build. The bibliography is
likewise generated: every DOI is fetched from Crossref at build time and an
unresolvable DOI fails the build.
